\documentclass{article}

% if you need to pass options to natbib, use, e.g.:
%     \PassOptionsToPackage{numbers, compress}{natbib}
% before loading neurips_2026

% The authors should use one of these tracks.
% Before accepting by the NeurIPS conference, select one of the options below.
% 0. "default" for submission
\usepackage{neurips_2026}
% the "default" option is equal to the "main" option, which is used for the Main Track with double-blind reviewing.
% 1. "main" option is used for the Main Track
%  \usepackage[main]{neurips_2026}
% 2. "position" option is used for the Position Paper Track
%  \usepackage[position]{neurips_2026}
% 3. "eandd" option is used for the Evaluations & Datasets Track
 % \usepackage[eandd]{neurips_2026}
% 4. "creativeai" option is used for the Creative AI Track
%  \usepackage[creativeai]{neurips_2026}
% 5. "sglblindworkshop" option is used for the Workshop with single-blind reviewing
 % \usepackage[sglblindworkshop]{neurips_2026}
% 6. "dblblindworkshop" option is used for the Workshop with double-blind reviewing
%  \usepackage[dblblindworkshop]{neurips_2026}
% 7. "education" option is used for the Education Track
%  \usepackage[education]{neurips_2026}


% After being accepted, the authors should add "final" behind the track to compile a camera-ready version.
% 1. Main Track
 % \usepackage[main, final]{neurips_2026}
% 2. Position Paper Track
%  \usepackage[position, final]{neurips_2026}
% 3. Evaluations & Datasets Track
 % \usepackage[eandd, final]{neurips_2026}
% 4. Creative AI Track
%  \usepackage[creativeai, final]{neurips_2026}
% 5. Workshop with single-blind reviewing
%  \usepackage[sglblindworkshop, final]{neurips_2026}
% 6. Workshop with double-blind reviewing
%  \usepackage[dblblindworkshop, final]{neurips_2026}
% Note. For the workshop paper template, both \title{} and \workshoptitle{} are required, with the former indicating the paper title shown in the title and the latter indicating the workshop title displayed in the footnote.
% For workshops (5., 6.), the authors should add the name of the workshop, "\workshoptitle" command is used to set the workshop title.
% \workshoptitle{WORKSHOP TITLE}
% 7. Education Track
% \usepackage[education, final]{neurips_2026}

% "preprint" option is used for arXiv or other preprint submissions
 % \usepackage[preprint]{neurips_2026}

% to avoid loading the natbib package, add option nonatbib:
%    \usepackage[nonatbib]{neurips_2026}

\usepackage[utf8]{inputenc} % allow utf-8 input
\usepackage[T1]{fontenc}    % use 8-bit T1 fonts
\usepackage{hyperref}       % hyperlinks
\usepackage{url}            % simple URL typesetting
\usepackage{booktabs}       % professional-quality tables
\usepackage{amsfonts}       % blackboard math symbols
\usepackage{nicefrac}       % compact symbols for 1/2, etc.
\usepackage{microtype}      % microtypography
\usepackage{xcolor}         % colors
\usepackage{amsmath,amssymb,amsthm}
\usepackage{graphicx}
\graphicspath{{../figures/}}
\newtheorem{proposition}{Proposition}
\newtheorem{definition}{Definition}

% Note. For the workshop paper template, both \title{} and \workshoptitle{} are required, with the former indicating the paper title shown in the title and the latter indicating the workshop title displayed in the footnote. 
\title{How Much Targeted Training Does It Take to Top a Multi-Task Leaderboard? Minimum Investment Under Different Aggregation Rules}

\author{Anonymous Authors}

\begin{document}
\maketitle

\begin{abstract}
A leaderboard turns per-task scores into one rank, so a developer can try to lift a chosen model by training on data related to individual sub-datasets. How much such targeted training is needed to put the model first, and how does that depend on the aggregation rule? We define the investment as the minimum raise of the target's task scores, within a per-task cap, that makes it strictly first, and measure it in points, tasks touched and largest single-task jump. We solve it exactly for the mean, median, minimum, Borda and Condorcet-winner rules, prove closed forms for the first three, and apply it to every model of the Open LLM Leaderboard (4495 targets, 6 tasks) and of HELM MMLU (78 targets, 57 subjects). With caps taken from observed fine-tuning gains, first place is reachable for only 12\% of OLLB targets under the mean (largest feasible rank 901) and 19\% of HELM targets (largest feasible rank 16); with a 100-point cap every target can reach it. Among feasible targets the mean is the costliest rule and the median and minimum the cheapest: at OLLB ranks 11--50 the other rules need 0.18--0.50 of the mean's total points. Which rule leaves most targets able to win is not stable: it changes with the cap and the dataset (on OLLB the mean at the data-grounded cap but the minimum at a uniform $+20$ points; on HELM the median at every cap tried). Conclusions depend on the cap, which on HELM is transferred from OLLB, and hold for fixed rivals and score tables only.
\end{abstract}

\section{Introduction}
\label{sec:intro}
Multi-task benchmarks rank models by aggregating per-task scores into one number, and the resulting rank influences which models are adopted, published and funded. A developer therefore has an incentive to improve a chosen target model specifically on the benchmark's sub-datasets, for instance by training on data related to them. How much such targeted improvement is needed to reach first place, and how that amount depends on the aggregation rule, is the question we study on public leaderboard data.

Our approach is deliberately simple and exact. We take the score table of a public leaderboard as fixed, let the target raise its own scores task by task up to a cap, and compute the \emph{minimum} raise that makes it strictly first under a given aggregation rule. The cap is the central modelling assumption: it stands for how far training on related data can move one score. We therefore report results under a data-grounded cap and an oracle cap, and report infeasibility (no allowed raise makes the target first) as an outcome rather than a missing value.

\paragraph{Contributions.}
(i) A formulation of the minimum targeted investment with three cost measures ($L_1$, $K$, $L_\infty$) and an explicit family of caps (Section~\ref{sec:setup}). (ii) Exact solvers for score rules (mean, median, minimum, power and trimmed means) and rank rules (Borda, Condorcet winner, plurality, Copeland), tested against brute force, with closed forms and proofs for the mean, minimum and median (Section~\ref{sec:solvers}, Appendix~\ref{app:proofs}). (iii) A study of every model on two public leaderboards, with bootstrap and evaluation-noise intervals (Sections~\ref{sec:exp}--\ref{sec:results}).

\paragraph{Findings.}
Under caps grounded in observed fine-tuning gains, first place is out of reach for most models, under every rule: the feasible fraction is 12\% (mean) on OLLB and 19\% on HELM. Where it is reachable, the mean rule costs the most points at every cap, and order-statistic and rank rules cost between 18\% and 50\% of that at OLLB ranks 11--50. Which rule leaves the most targets able to win, in contrast, depends on the cap and on the dataset, so no rule is uniformly the most exposed. Two of three hypotheses fixed beforehand are only partly borne out (Section~\ref{sec:results}). All conclusions are conditional on the cap and on the other models being fixed; Section~\ref{sec:disc} lists where they fail.

\section{Related work}
\label{sec:related}
\paragraph{Benchmarks and leaderboards.}
MMLU \citep{hendrycks2021mmlu} aggregates 57 subjects, and HELM \citep{liang2023helm} releases per-scenario scores and raw model outputs for many models; BIG-Bench Hard \citep{suzgun2023bbh} selects hard tasks from BIG-Bench. These provide the multi-task score matrices that we analyse.

\paragraph{Fragility of rankings.}
\citet{dehghani2021lottery} show that the relative order of methods can change with the choice of benchmark tasks (the ``benchmark lottery''). \citet{singh2025illusion} document how private testing and unequal data access distort the Chatbot Arena leaderboard. Both concern what is measured and disclosed; we instead hold the leaderboard fixed and ask what a single developer must change for one model to win.

\paragraph{Aggregation and social choice.}
\citet{lanctot2025vase} treat each task as a voter and use voting theory to build cross-task evaluation; the Handbook of Computational Social Choice \citep{brandt2016handbook} covers the rules (Borda, Condorcet-consistent rules, plurality) we reuse. Our rules are the same families, but our question is the cost of changing the winner, not the choice of rule.

\paragraph{Training on the test task and contamination.}
\citet{zhou2023cheater} show that benchmark leakage into training can raise scores sharply, and \citet{yang2023rephrased} show that rephrased test samples evade common decontamination. \citet{chen2026incentives} analyse, in a game-theoretic model, the incentives that benchmarks create for strategic post-training. Our analysis is empirical and combinatorial: given a fixed score table, we compute the exact minimum targeted improvement per rule, under explicit assumptions about how much any one task can be improved.

\paragraph{Positioning.}
Closest in question is \citet{chen2026incentives}, who model developers choosing how much to spend on benchmark-specific improvement of a single benchmark score, prove non-existence of equilibria for current protocols and a unique equilibrium for tune-before-test. That analysis concerns one score per model; the choice of aggregation rule over tasks does not enter it. We take the opposite stance: no game and no equilibrium, a fixed table, and the exact cost for one developer under each of several aggregation rules. \citet{colombo2022best} propose a social-choice-based aggregation for NLP benchmarks and show it changes conclusions relative to the mean; we do not propose a rule but measure how costly it is to move a model to first place under each. The computational-social-choice literature studies the analogous \emph{bribery} problem, the cost of changing votes so that a chosen candidate wins under a voting rule \citep{faliszewski2009bribery,brandt2016handbook}. Our setting differs in that the manipulator changes its own scores on tasks (the ``voters'') and pays in score points under per-task caps, which gives closed forms for score rules and integer programs for rank rules.

\section{Problem setup}
\label{sec:setup}
\subsection{Leaderboard, target and investment}
A leaderboard is a score matrix $X\in[0,100]^{M\times n}$ of $M$ models on $n$ tasks. Fix a target model $i$ with score vector $x=X_{i\cdot}$; all other rows stay fixed. In targeted training the developer can raise the target's score on task $t$ to any $y_t\in[x_t,c_t]$, where $c_t$ is a \emph{cap} (Section~\ref{sec:caps}). An aggregation rule $f$ maps the score vectors to a ranking; the target is \emph{first} if it is strictly above every other model under $f$ (by a margin $\varepsilon=10^{-3}$ points for score rules). We report the investment in three ways:
\[
L_1=\sum_t (y_t-x_t),\qquad K=\#\{t: y_t>x_t\},\qquad L_\infty=\max_t (y_t-x_t),
\]
the total points, the number of tasks touched, and the largest single-task jump. The minimum of each is taken separately; $L_\infty$ is also the cost of the uniform ``general improvement'' that raises every task by the same amount. For comparison across rules and datasets we also report $K/n$ and $L_1/n$.

\subsection{Caps: how much can one task be improved?}
\label{sec:caps}
Without a bound on $y_t$ the question is trivial (put every task at 100). The cap $c_t$ is therefore the central modelling assumption, and we never collapse it to a single choice. Two caps appear in the main text: \texttt{ceil}, $c_t=100$, an oracle upper bound on what can be done and so a lower bound on any required investment; and \texttt{emp99}, $c_t=\min(100,x_t+g_t^{99})$ where $g_t^{99}$ is the 99th percentile of the observed gains (fine-tune minus its declared base model) on task $t$ in the Open LLM Leaderboard. \texttt{emp99} is the only data-grounded cap. It inherits a selection bias, since only published fine-tunes appear, and it measures gains attained by the community, not the gains a determined developer could attain. Further caps (best existing score $+0$ or $+2$, and at most $+5/+10/+20$ points per task) are reported in the appendix. When no $y\le c$ makes the target first, the case is \emph{infeasible}; infeasibility is reported as a rate and never coded as a large investment, and all investment statistics are conditional on feasibility.

\subsection{Aggregation rules}
\label{sec:rules}
Score rules apply a function $f$ to the score vector of each model: the \emph{mean} (the rule of both leaderboards we use), the \emph{median}, and the \emph{minimum} (worst-task score). Rank rules use only the order of models on each task: \emph{Borda} (sum over tasks of the number of models beaten, ties counting one half; equivalent to mean rank) and the \emph{Condorcet-winner} requirement (the target beats every rival on a strict majority of tasks, ties counting one half). Being a Condorcet winner is sufficient for being first under any Condorcet-consistent rule such as Copeland or Schulze, so its investment upper-bounds theirs. Geometric, harmonic and power means, the trimmed mean and plurality are in the appendix.

\subsection{Exact minimum-investment solvers}
\label{sec:solvers}
For a score rule let $\tau=\max_{j\neq i} f(X_{j\cdot})$ be the level of the best rival; the target is first iff $f(y)\ge \tau+\varepsilon$. For the mean, the minimum and the median the minimum investment has a closed form (Propositions~\ref{prop:mean}--\ref{prop:median}, proved in Appendix~\ref{app:proofs}); for power means and the trimmed mean we solve a convex program or enumerate exactly. For rank rules a raise of task $t$ to level $s+\varepsilon$ passes exactly the models with score at most $s$ on that task, so only these finitely many levels need be considered, and we solve the resulting integer program exactly with HiGHS, adding rival constraints lazily. Every solver is checked against brute-force search on small random instances (unit tests in the supplementary code). For each target we report $L_1$ and $K$, each minimised separately, and $L_\infty$ for the uniform baseline.

\begin{proposition}[Mean]\label{prop:mean}
Let $f$ be the mean and $h_t=c_t-x_t$. The target can become first iff $\sum_t c_t> n\tau$ (up to $\varepsilon$). The minimum $L_1$ is $n(\tau-f(x))$ regardless of how it is spread, the minimum $K$ is the least $k$ for which the $k$ largest $h_t$ sum to at least this amount, and the minimum $L_\infty$ is $\tau-f(x)$ (raise all tasks equally) when no cap binds.
\end{proposition}
\begin{proposition}[Minimum]\label{prop:min}
Let $f=\min_t y_t$. The target can become first iff $c_t>\tau$ for all $t$ with $x_t\le\tau$. Then $L_1=\sum_t(\tau-x_t)^+$ and $K=\#\{t:x_t\le\tau\}$: every deficient task must be raised.
\end{proposition}
\begin{proposition}[Median, $n$ odd]\label{prop:median}
Let $m=(n+1)/2$, $d_t=(\tau-x_t)^+$ for tasks with $c_t>\tau$ and $d_t=\infty$ otherwise. The target is first iff at least $m$ tasks end above $\tau$. The minimum $L_1$ is the sum of the $(m-\#\{t:x_t>\tau\})^+$ smallest finite $d_t$ among tasks with $x_t\le\tau$, and $K$ equals that count; the target can be made first iff enough finite $d_t$ exist.
\end{proposition}
(All statements are for $\varepsilon\to0$; the code uses $\varepsilon=10^{-3}$ points. For even $n$ the median is the average of the two middle order statistics and we only have the computed solution.) We verify Propositions~\ref{prop:mean}--\ref{prop:median} numerically against the solvers (Section~\ref{sec:exp}); no theorem is claimed for the rank rules.

\section{Experimental setup}
\label{sec:exp}
\paragraph{Data.}
We use two public leaderboards, each at a fixed, recorded version (checksums in the supplementary reproduction notes). \emph{HELM MMLU}, release v1.13.0 (summary dated 2025-01-10): the exact-match score of 79 models on each of the 57 MMLU subjects ($79\times57$, no missing cell). \emph{Open LLM Leaderboard} (OLLB) \texttt{contents}, dataset commit \texttt{9c09a7c} (2025-03-20): 4576 rows, of which 4497 have distinct model names after we keep the best-scoring variant of each duplicated name, and 4496 after removing one flagged model; each model has six normalised task scores (IFEval, BBH, MATH Lvl 5, GPQA, MUSR, MMLU-PRO), so $n=6$. Merged models are kept in the main analysis and dropped in a robustness slice. The mean we use equals the sites' own aggregate (OLLB \texttt{Average}, HELM ``MMLU All Subjects'') to within $6\times10^{-14}$ points on every model. Neither source declares a licence for the score data; we redistribute only the derived score tables.

\paragraph{Targets, caps and what is reported.}
Every model other than the current leader is a target in turn, with all others fixed (4495 targets on OLLB, 78 on HELM). On OLLB the primary cap is \texttt{emp99} (per-task 99th percentile of the gain of a fine-tune over its declared base model, estimated from 540 such pairs in the same leaderboard). HELM has no base/fine-tune pairs, so its primary cap, \texttt{emp99T}, applies the OLLB MMLU-PRO 99th-percentile gain to all 57 subjects; this is a transfer assumption, not a measurement on HELM. The oracle cap \texttt{ceil} is reported beside it. All investment statistics are medians over feasible targets, shown together with the feasible fraction. OLLB rank rules are exact integer programs, which at $M=4496$ models we run only for the 50 best targets under the mean, and, as a separate slice, for every target on the 500 best models taken as a leaderboard of their own; we apply a 120\,s limit per instance and report time-outs as censored, never as infeasible (one instance on the slice). Robustness runs (evaluation noise, task bootstrap, margin, slices, weights, cost scale) are described with their results.

\section{Results}
\label{sec:results}
All statements below are about targets that can be made first within the cap; the feasible fraction is shown beside every cost, and a rule is called cheaper than another only on targets that are feasible under both.

\paragraph{The closed forms are exact (Q1).}
Propositions~\ref{prop:mean}--\ref{prop:median} agree with the exact solvers on every target of both datasets and every cap: the largest deviation in $L_1$ is below $2\times10^{-12}$ points, the deviation in $K$ is zero, and no feasibility verdict differs (Appendix Table~\ref{tab:T2}). The OLLB median is solver-only because $n=6$ is even.

\begin{figure}[t]
\centering
\includegraphics[width=\linewidth]{F_ollb_primary.pdf}\\[2pt]
\includegraphics[width=\linewidth]{F_helm_primary.pdf}
\caption{Minimum targeted investment to become first versus the target's current rank under the mean, primary caps. Left: fraction of targets that can reach first at all. Middle and right: tasks touched ($K/n$) and total points per task ($L_1/n$), \emph{among feasible targets}. Top: OLLB, $n=6$, cap \texttt{emp99} (4495 targets; Borda and Condorcet only for the 50 best). Bottom: HELM MMLU, $n=57$, cap \texttt{emp99T} (78 targets; the cap is transferred from OLLB). Points are bin medians (bins double from rank 2; bins with fewer than 3 targets omitted), bands are 95\% bootstrap intervals over targets. Generated by \texttt{figures/src/fig\_rank\_rules.py}.}
\label{fig:main}
\end{figure}

\paragraph{Feasibility depends on rank and on the rule (Q3).}
Under the data-grounded caps, targeted training cannot make most models first (Figure~\ref{fig:main}, left). On OLLB the feasible fraction under the mean and the median is close to one for targets up to rank a few hundred and near zero beyond rank about one thousand; on HELM every target at ranks 2--10 is feasible under the mean, median, Borda and Condorcet rules, and the fraction falls steeply afterwards (Table~\ref{tab:T3}: for ranks 11--50 only 6 of 40 targets are feasible under the mean and 17 of 40 under the median). The minimum rule is the most restrictive on both datasets, since every task below the best rival's weakest score must be raised and any one cap can block it (Proposition~\ref{prop:min}). The cap, not the rule, decides feasibility far from the top: with the oracle cap \texttt{ceil} every target is feasible (Appendix Figure~\ref{fig:oracle}).

Table~\ref{tab:T4} gives the frontier. On OLLB the fraction of targets that can reach first place under \texttt{emp99} is 12\% under the mean, 10\% under the median and 6\% under the minimum; every target is feasible only up to rank 394 (mean), 90 (median) and 5 (minimum), and the largest feasible rank is 901, 1207 and 2225. On HELM under \texttt{emp99T} the fractions are 19\%, 33\% and 12\%, every target is feasible up to rank 16, 19 and 6, and the largest feasible rank is 16, 29 and 20; the Condorcet requirement reaches rank 26 and Borda 17. On HELM the median is more exposed to targeted training than the mean (33\% against 19\% of targets feasible); on OLLB it is slightly less exposed by share (10\% against 12\%, non-overlapping intervals) although its largest feasible rank is deeper, so the comparison of median and mean is not the same on the two datasets. The minimum is feasible rarely but, when feasible, cheaply. The largest feasible rank can be an isolated case (OLLB, minimum), which is why the column ``all feasible up to'' is shown.

\begin{table}[t]
\caption{Feasibility frontier under the primary cap (\texttt{emp99} on OLLB, \texttt{emp99T} on HELM) and the oracle cap (\texttt{ceil}). \emph{share}: fraction of targets for which first place is reachable (Wilson 95\% intervals); $r^*$: largest rank under the mean among feasible targets (may be an isolated case); \emph{all feas.\ up to}: largest $r$ such that every target of rank $\le r$ is feasible; \emph{first bin $<$ 1/2}: first rank bin (bins double from rank 2) in which fewer than half are feasible. $^\dagger$: the 50 best targets under the mean only, not comparable with the score rules' fractions. Generated by \texttt{figures/src/tables.py}.}
\label{tab:T4}
\centering\scriptsize
\resizebox{\linewidth}{!}{\input{../tables/T4_frontier_ollb.tex}}\\[4pt]
\resizebox{\linewidth}{!}{\input{../tables/T4_frontier_helm.tex}}
\end{table}

\paragraph{The rule that is easiest to win under depends on the cap.}
Changing the cap changes which rule leaves most targets feasible (Appendix Table~\ref{tab:T5}, Figures~\ref{fig:capgrid} and~\ref{fig:capgrid_helm}). On OLLB, with uniform caps of $+5$, $+10$ and $+20$ points per task, the feasible share is $0.3\%$, $2.5\%$ and $18.5\%$ under the mean, $0.9\%$, $4.5\%$ and $14.3\%$ under the median and $2.3\%$, $9.9\%$ and $63.9\%$ under the minimum; with the per-task data-grounded cap \texttt{emp99} it is $12.1\%$, $10.0\%$ and $6.1\%$. The minimum thus goes from the least to the most exposed rule as the caps become uniform and loose, which we read as an effect of the per-task caps (the minimum needs every deficient task raised, Proposition~\ref{prop:min}, so one tight task cap blocks it); we have not isolated the individual task caps. On HELM the median is at least as exposed as the mean at every cap tried ($20.5\%$ against $14.1\%$ at $+5$, $51.3\%$ against $37.2\%$ at $+10$, $84.6\%$ against $71.8\%$ at $+20$, $33.3\%$ against $19.2\%$ at \texttt{emp99T}), and the Condorcet requirement is at least as exposed as the mean at all of them (e.g.\ $32.1\%$ against $19.2\%$ at \texttt{emp99T}; both reach $98.7\%$ at \texttt{emp99P}). On the 500 best OLLB models taken as a leaderboard of their own, where rank rules are exact for every target, the Condorcet requirement is the least exposed rule at \texttt{emp99}, the reverse of HELM: the feasible share is $94.8\%$ under the mean, $90.6\%$ under Borda, $79.4\%$ under the median, $56.3\%$ under plurality, $40.3\%$ under the minimum and $30.5\%$ under the Condorcet requirement (Table~\ref{tab:T6}). The requirement is to win a strict majority of tasks against every rival (four of six here, 29 of 57 on HELM). So the claim ``rule $A$ is easier to manipulate than rule $B$'' has to be stated with a dataset and a cap. What is stable is the cost among targets feasible under both rules: the median and minimum rules (and Borda and Condorcet, where solved) cost less than the mean at every cap, with median ratios of $L_1$ to the mean between $0.08$ and $0.61$ (median), $0.09$ and $0.35$ (minimum) on OLLB and between $0.04$ and $0.18$ (median), $0.05$ and $0.15$ (minimum) on HELM (\texttt{results/summary/paired\_*.csv}; the number of jointly feasible targets is small at tight caps, as few as 5 at $+5$ points). On HELM the exact Copeland solution coincides with the Condorcet requirement (same feasible set, $L_1$ equal up to the margin $\varepsilon$ at \texttt{emp99T} and \texttt{ceil}), so the Condorcet requirement is tight there.

\paragraph{Cost, conditional on feasibility, differs strongly by rule (Q2).}
Among feasible targets, the mean rule is the most expensive on both datasets and the minimum and median rules the cheapest (Figure~\ref{fig:main}, right; Table~\ref{tab:T3}). On OLLB at ranks 11--50 the median feasible target needs $37.8$ points in total under the mean, $9.1$ under the median, $6.9$ under the minimum, $11.7$ under Borda and $20.5$ under the Condorcet-winner requirement; relative to the mean on jointly feasible targets, the ratios are $0.22$, $0.18$, $0.27$ and $0.50$. On HELM at ranks 11--50 the corresponding totals are $321$, $38$, $10$, $188$ and $105$ points. The reason is structural: the mean must close the whole summed gap ($n$ times the gap, Proposition~\ref{prop:mean}), whereas order-statistic rules pay only for the cheapest majority of tasks or the weakest tasks. Rank rules sit in between, because they pay for passing rivals task by task. Table~\ref{tab:T3} shows intervals; the ordering mean $>$ Condorcet $>$ Borda $\gtrsim$ median $\gtrsim$ min on OLLB is not uniform across rank bins (the intervals at ranks 2--10 are wide, with $n=9$ targets).

\begin{table}[t]
\caption{Median investment to become first among feasible targets, by rank stratum and rule, primary caps: OLLB (\texttt{emp99}, top) and HELM (\texttt{emp99T}, bottom). ``feas.'' = feasible targets / targets in the stratum; $K$, $L_1$ with 95\% bootstrap intervals over targets; the last column is the median ratio of $L_1$ to the mean rule over targets feasible under both. OLLB Borda and Condorcet use the 50 best targets only. Generated by \texttt{figures/src/tables.py}.}
\label{tab:T3}
\centering\scriptsize
\resizebox{\linewidth}{!}{\input{../tables/T3_rules_ollb_emp99.tex}}\\[4pt]
\resizebox{\linewidth}{!}{\input{../tables/T3_rules_helm_emp99T.tex}}
\end{table}

\paragraph{Which of our predictions held.}
The plan fixed three hypotheses before the full results (H1--H3, Appendix~\ref{app:hyp}). H2 (cost rises with rank, feasibility falls as caps tighten, rule ordering can change with the cap) holds on both datasets, including the last clause, as shown above. H1 (order-statistic and majority rules need \emph{more} tasks but possibly fewer points than the mean) holds only in part: they need fewer points, but not more tasks, except that the Condorcet requirement touches more tasks than the mean on OLLB at $n=6$ ($K=4$ against $2$ at ranks 11--50); on HELM the median needs fewer tasks than the mean ($K=12$ against $46$). H3 is examined next.

\paragraph{Number of tasks (Q4).}
Figure~\ref{fig:scaling} draws random subsets of $n\in\{3,6,12,24,48\}$ HELM subjects and takes the model ranked second or fifth within each subset as the target. For the rank-2 target $K/n$ falls with $n$ under every rule, and for the rank-5 target the mean and Borda rules keep touching about $0.4$ of the tasks while the median, minimum and plurality fall. For the rank-5 target the mean rule's $L_1/n$ stays roughly constant in $n$ while that of the median, minimum and plurality falls; for the rank-2 target all rules fall. The first half of H3 (that $K/n$ of the median and majority rules does not shrink) is therefore not supported for the median. The second half (that the mean rule's $L_1/n$ depends on the density of gaps) was not tested: random subsets of HELM subjects do not vary the density independently, and we only observe that $L_1/n$ is roughly constant in $n$ for the rank-5 target. The bands are wide (60 subset draws), so we make no claim about differences between neighbouring $n$. OLLB, with $n=6$, is a single small-$n$ point and supports no scaling claim.

\begin{figure}[t]
\centering
\includegraphics[width=\linewidth]{F_helm_scaling.pdf}
\caption{HELM: investment versus the number of tasks $n$ in a random subset of subjects, for the target ranked second (left) or fifth (right) within the subset under the mean; bands span the 2.5--97.5 percentiles over 60 subset draws, among feasible draws, cap \texttt{emp99T}. Generated by \texttt{figures/src/fig\_scaling.py}.}
\label{fig:scaling}
\end{figure}

\paragraph{Robustness (Q5).}
We resample evaluation noise (binomial standard errors from item counts, Gaussian draws; 300 draws for score rules on OLLB, 60 for its rank rules, 200 for all rules on HELM) and tasks (bootstrap over tasks with each task's cap following it; same draw counts) for seven OLLB and six HELM targets spread over ranks 2--100 (Figure~\ref{fig:robust}, Appendix). For every cell in which the observed data are feasible (26 HELM cells under each of the two resampling schemes, 24 OLLB score-rule cells and 16 OLLB rank-rule cells), the observed $L_1$ lies inside the 2.5--97.5 percentile range of the feasible draws; it lies below the draw median in 18 of 26 (HELM noise), 25 of 26 (HELM bootstrap), 23 of 24 and 20 of 24 (OLLB score rules, noise and bootstrap) and 14 of 16 and 15 of 16 (OLLB rank rules). The observed cost is therefore a typical to slightly optimistic value of what resampled leaderboards require, with a median relative gap to the draw median of $8$--$27\%$ over the six groups. Feasibility is the fragile part: the HELM rank-20 target is infeasible on the observed data under the mean but feasible in $4.5\%$ of noise draws ($0.88$ infeasible under the median, $0.70$ under the minimum), whereas the rank-40 target is infeasible in every draw under the mean and median and in at least $84\%$ of draws under every rule; under the task bootstrap the share of infeasible draws is $0.22$--$0.87$ for the OLLB Condorcet requirement and $0.35$--$0.92$ for plurality. The intervals for the median and minimum are wide, and a large share of resampled leaderboards make the minimum rule infeasible. A target whose gap to first place is within the evaluation noise therefore cannot be said to need any particular investment: a win by the margin $\varepsilon$ would not survive resampling.

\paragraph{Sensitivity to modelling choices.}
Appendix Table~\ref{tab:T6} compares variants with the baseline. A margin of 1 or 2 points in place of $10^{-3}$ lowers the feasible share (HELM mean $19.2\%\to15.4\%$ and $11.5\%$, OLLB mean $12.1\%\to11.1\%$ and $10.4\%$) and raises the cost on targets feasible in both (median ratio $1.08$--$1.73$ for the mean, $1.15$--$2.34$ for the median) without changing the ordering of the rules by cost. Dropping the merged models (3788 targets) leaves the OLLB mean share at $11.5\%$ and the medians of $L_1$ at $74.1$ (mean), $26.6$ (median) and $7.4$ (minimum) points. A cost measured in standard deviations of each task instead of points keeps the order mean $>$ median $>$ minimum on HELM (median $L_1$ in sigma units at the primary cap: $17.5$, $2.5$, $0.6$) and mean $>$ minimum on OLLB ($4.3$ and $1.6$; the median is not computed there because $n=6$ is even). Scoring OLLB in raw accuracy instead of normalised points keeps the order mean $>$ median $>$ minimum at the oracle cap (median $L_1$ of $154.1$, $92.9$ and $74.7$ points), but the minimum is then twice as expensive relative to the mean ($0.48$ against $0.24$ of the baseline), and a uniform $+10$ cap ranks the rules differently (feasible shares $6.6\%$, $0.5\%$, $5.0\%$ for mean, median and minimum on raw scores against $2.5\%$, $4.5\%$, $9.9\%$ on normalised scores): uniform caps are not scale-free. We use the normalised score because the leaderboard computes its average on it. Random positive task weights (Dirichlet$(1)$, Dirichlet$(5)$; 500 draws for each stratified target) lower the investment relative to the uniform mean (median ratio $0.54$ and $0.70$ on OLLB, $0.43$ and $0.68$ on HELM) since a point raised on a heavily weighted task counts for more; they define a different rule rather than checking the mean, and we do not decompose the reduction.

\section{Discussion and limitations}
\label{sec:disc}
\paragraph{What the results say.}
Under caps taken from observed fine-tuning gains, a developer who can lift only the target's own scores, and whose rivals stand still, tops these leaderboards only from near the top. Two statements held for every cap, dataset and variant we ran. First, where first place is reachable the mean rule demands the most points: it must close $n$ times the gap (Proposition~\ref{prop:mean}) whatever the spread, while the median and minimum pay only for a majority of tasks or for the weak tasks and the rank rules for passing rivals task by task. Second, under the mean the cost of a feasible target does not depend on the cap; the cap decides only whether the target can get there. By the largest feasible rank $r^*$ the mean is the lowest of the three score rules on both datasets (OLLB $901$, $1207$, $2225$ for mean, median, minimum; HELM $16$, $29$, $20$), but $r^*$ can be an isolated case, and by the share of feasible targets and the fully feasible prefix the minimum is the lowest. A third, tempting statement does not hold: that the mean is also the hardest rule to win under in the sense of reachability. On HELM the median and the Condorcet requirement are at least as exposed as the mean at every cap and Borda reaches only one rank further than the mean (17 against 16); on OLLB the mean is the most exposed score rule at the data-grounded cap, the minimum the most exposed at uniform $+20$ points, and on the 500-model slice the Condorcet requirement is the least exposed rule. Cheap in points and reachable are different properties of a rule, and only the first is stable. These orderings come with caveats we state next.

\paragraph{Conditions under which the conclusions fail or are not established.}
(1) \emph{The cap.} Feasibility is determined by the cap, not by the rule, far from the top (oracle rows of Table~\ref{tab:T4}). The OLLB cap uses gains of published fine-tunes over their declared base models; this is selection-biased toward successes and says nothing about what a determined developer could achieve with a new data source. The HELM cap is a transfer of the OLLB MMLU-PRO gain to 57 subjects and has not been measured on HELM. Headline claims are therefore stated per dataset. (2) \emph{Independence across tasks.} We let each task be raised at no cost to others; targeted training can trade off tasks, which would make the true cost higher. (3) \emph{Fixed rivals.} Other developers do not respond; leaderboards with many strategic developers need a game model such as that of \citet{chen2026incentives}. (4) \emph{Score tables, not training.} The investment is in score points, not compute, data or money, so ``cheap'' means few score points, not little effort. (5) \emph{Ranks within noise.} The margin of $10^{-3}$ points that defines ``first'' is below the evaluation noise of the scores; a win by that margin would not survive resampling (Figure~\ref{fig:robust}), and the intervals for the median and minimum are wide. (6) \emph{Rule coverage and size.} On the full OLLB board rank rules were solved exactly for the 50 best targets only, so their feasibility fractions there are not comparable with those of the score rules (the 500-model slice, Table~\ref{tab:T6}, is a different leaderboard with different ranks); the Condorcet requirement is a sufficient condition for Condorcet-consistent rules and therefore an upper bound for Copeland-type rules, tight on HELM where Copeland was solved exactly; Kemeny and Schulze were not studied. (7) \emph{Small $n$ and small $M$.} OLLB has six tasks and HELM 78 targets, so rank bins on HELM hold few targets and their intervals are wide; scaling in $n$ comes from HELM subsets only and is a statement about random subsets of one benchmark family. (8) \emph{Data snapshot.} Both leaderboards are frozen snapshots (OLLB 2025-03, HELM v1.13.0); we do not claim anything about later models. (9) \emph{Unfinished items.} Rank rules are exact for the 50 best OLLB targets on the full board and for all targets only on the 500-model slice (where one instance hit the time limit); Copeland was solved on HELM only; an extension of OLLB to its roughly 38 subtasks (more tasks at the same leaderboard) was not run; evaluation-noise and bootstrap intervals cover seven OLLB and six HELM targets per rule, not all targets, and use 60 draws for OLLB rank rules.

\paragraph{Broader impact.}
The analysis tells leaderboard designers how exposed each rule is to targeted training under stated assumptions, which supports choosing rules and caps (for instance holding out sub-datasets) deliberately. It could also be read as a recipe for gaming a ranking. We think the risk is small: it needs no new capability beyond fine-tuning on data a developer already can use, and our results indicate that for most models it is infeasible under the observed gains. We do not evaluate any model's training.

\section{Reproducibility}
\label{sec:repro}
All numbers, tables and figures are produced by the supplementary code from the pinned public data (versions, checksums, seeds, commands, expected run times and a table mapping every figure and table to its generating script are in \texttt{REPRODUCE.md}). The solvers are deterministic; random draws use seed 0. The full pipeline needs no GPU and takes roughly eight hours of wall-clock on a shared container with two CPUs and 4\,GB (Xeon Gold 6230; per-step runtimes in \texttt{REPRODUCE.md}). Neither leaderboard declares a licence for its score data; we release only derived score tables. No material about the published study that motivated this question was searched for or read, and none was encountered (reproduction notes, section on forbidden material).


\bibliographystyle{plainnat}
\bibliography{refs}

\appendix
\section{Proofs of Propositions~\ref{prop:mean}--\ref{prop:median}}
\label{app:proofs}
\paragraph{Proposition~\ref{prop:mean}.} $f(y)=\frac1n\sum_t y_t\ge\tau$ iff $\sum_t(y_t-x_t)\ge n(\tau-f(x))$, and $\sum_t(y_t-x_t)\le\sum_t h_t$ with each $0\le y_t-x_t\le h_t$; so feasibility is $\sum_t h_t\ge n(\tau-f(x))$, i.e. $\sum_t c_t\ge n\tau$, and the minimum of the (linear) total is the requirement itself, attained by any feasible spread. $K$ is minimised by filling the largest headrooms first, since touching $k$ tasks can supply at most the sum of the $k$ largest $h_t$. The uniform raise by $\tau-f(x)$ attains $L_\infty$, and no spread has a smaller maximum because the average raise is $\tau-f(x)$.
\paragraph{Proposition~\ref{prop:min}.} $\min_t y_t>\tau$ requires $y_t>\tau$ for every $t$, hence $c_t>\tau$ wherever $x_t\le\tau$ (tasks with $x_t>\tau$ need no change). The cheapest $y_t$ on a deficient task is $\tau$ (up to $\varepsilon$), so $L_1=\sum_t(\tau-x_t)^+$, and $K$ counts the deficient tasks. Both are forced.
\paragraph{Proposition~\ref{prop:median}.} For odd $n$ the median exceeds $\tau$ iff at least $m$ coordinates exceed $\tau$. Tasks with $x_t>\tau$ already count. We need $m'=(m-\#\{x_t>\tau\})^+$ more tasks. Lifting task $t$ above $\tau$ costs at least $d_t$ (and is possible iff $c_t>\tau$); costs are separable, so the cheapest set is the $m'$ smallest finite $d_t$ and $K=m'$, which cannot be smaller because each lifted task adds one to the count.


\section{Further results}
\begin{table}[h]
\caption{Closed forms (Propositions~\ref{prop:mean}--\ref{prop:median}) against the exact solvers and an independent generic LP/MILP (60-target sample), over all caps. Generated by \texttt{figures/src/tables.py} from \texttt{results/theory\_check.csv}.}
\label{tab:T2}
\centering\scriptsize
\input{../tables/T2_theory_check.tex}
\end{table}
\begin{table}[h]
\caption{Share (\%) of targets for which first place is reachable, by rule and per-task cap (columns roughly from tight to loose; $+h$: at most $h$ points per task; \texttt{emp95}, \texttt{emp99}: per-task percentile of observed fine-tune gains; $^T$: transferred from OLLB MMLU-PRO; pooled: pooled six-task percentile; front: best rival score; 100: score ceiling). Wilson 95\% intervals below each cell. $r^*$ and ``all feas.\ up to'' are at the primary cap. Top: OLLB; Borda, Condorcet and plurality were solved for the 50 best targets at \texttt{emp99} and the oracle cap only and are not comparable with the score rules' shares; empty cells were not computed. Bottom: HELM. Generated by \texttt{figures/src/tables.py} from \texttt{results/summary/frontier\_*.csv}.}
\label{tab:T5}
\centering\scriptsize
\resizebox{\linewidth}{!}{\input{../tables/T5_capshare_ollb.tex}}\\[4pt]
\resizebox{\linewidth}{!}{\input{../tables/T5_capshare_helm.tex}}
\end{table}
\begin{table}[h]
\caption{Sensitivity variants against the baseline: margin of 1 or 2 points instead of $10^{-3}$; OLLB without merged models; OLLB top-500 slice (499 targets, exact Borda, Condorcet and plurality); OLLB raw accuracies (the $+10$ cap is in raw points, so its baseline column is not a like-for-like comparison of cap size). Feasible share, median $L_1$ and $K$ among feasible targets, and the median per-target ratio of $L_1$ to the baseline over targets feasible in both. Generated by \texttt{figures/src/tables.py} from \texttt{results/summary/summary\_variants.csv}.}
\label{tab:T6}
\centering\scriptsize
\resizebox{\linewidth}{!}{\input{../tables/T6_variants.tex}}
\end{table}
\begin{figure}[h]
\centering
\includegraphics[width=\linewidth]{F_ollb_oracle.pdf}\\[2pt]
\includegraphics[width=\linewidth]{F_helm_oracle.pdf}
\caption{As Figure~\ref{fig:main} with the oracle cap \texttt{ceil} (scores up to 100): a lower bound on any required investment; every target is feasible, so the differences between rules are not caused by the cap.}
\label{fig:oracle}
\end{figure}
\begin{figure}[h]
\centering
\includegraphics[width=\linewidth]{F_ollb_capgrid.pdf}
\caption{OLLB: effect of the cap on feasibility (left) and cost among feasible targets (right) for three rank strata. Caps run from tight to loose (+5 points per task, \texttt{emp95}, +10, \texttt{emp99}, +20, best rival score, 100). For the mean rule the cost among feasible targets does not depend on the cap (Proposition~\ref{prop:mean}); only the feasible fraction does. Generated by \texttt{figures/src/fig\_cap\_grid.py}.}
\label{fig:capgrid}
\end{figure}

\begin{figure}[h]
\centering
\includegraphics[width=\linewidth]{F_helm_capgrid.pdf}
\caption{HELM: effect of the cap on feasibility (left) and cost among feasible targets (right) for three rank strata (4, 15 and 58 targets). Caps are ordered roughly from tight to loose; the order is approximate because uniform caps (+h points) and per-task percentile caps (\texttt{emp}) are not strictly comparable, and \texttt{emp99T} lies below $+10$ points. Bands: Wilson 95\% (left), 95\% bootstrap over targets (right). Among feasible targets the cost of the mean rule changes with the cap only because the set of feasible targets changes; the rank 2--5 stratum has 4 targets and its intervals carry no ranking of the rules. Generated by \texttt{figures/src/fig\_cap\_grid.py}.}
\label{fig:capgrid_helm}
\end{figure}
\begin{figure}[h]
\centering
\includegraphics[width=\linewidth]{F_robust.pdf}
\caption{Stability of the minimum total gain $L_1$ (top rows) and share of resampled leaderboards in which first place is unreachable (bottom rows) under evaluation noise (left) and a bootstrap over tasks (right), for OLLB (top, cap \texttt{emp99}, seven targets of mean-rank 2--100) and HELM (bottom, cap \texttt{emp99T}, six targets of rank 2--40). Dots and bars: median and 2.5--97.5 percentile range over feasible draws (a spread of draws, not a confidence interval of the median); open diamonds: observed data. Draws: 300 (OLLB score rules), 60 (OLLB Borda and Condorcet, exact integer programs), 200 (HELM, all rules). Where only the lower row shows a point the target is infeasible in (almost) every draw and on the observed data. Generated by \texttt{figures/src/fig\_robust.py}.}
\label{fig:robust}
\end{figure}

\section{Hypotheses fixed before the full results}
\label{app:hyp}
The plan fixed three hypotheses from the closed forms before the full result set was available. H1: with loose caps, score-sum rules need the same $L_1$ regardless of $K$, while order-statistic and majority rules need more tasks but possibly fewer points than the mean. H2: required investment grows with the target's initial rank; feasibility falls as caps tighten; rule ordering can change with the cap. H3: for fixed per-task caps the number of tasks that majority and median rules need does not shrink with $n$, while the mean rule's $L_1/n$ depends on the density of gaps. Anything beyond these is exploratory.


\newpage
\section*{NeurIPS Paper Checklist}




\begin{enumerate}

\item {\bf Claims}
    \item[] Question: Do the main claims made in the abstract and introduction accurately reflect the paper's contributions and scope?
    \item[] Answer: \answerYes{}
    \item[] Justification: Abstract and introduction state the two dataset-specific findings, the cap assumption and their scope; each is backed by Table~\ref{tab:T3}, Table~\ref{tab:T4} and Figure~\ref{fig:main}, and Section~\ref{sec:disc} lists where they fail.
    \item[] Guidelines:
    \begin{itemize}
        \item The answer \answerNA{} means that the abstract and introduction do not include the claims made in the paper.
        \item The abstract and/or introduction should clearly state the claims made, including the contributions made in the paper and important assumptions and limitations. A \answerNo{} or \answerNA{} answer to this question will not be perceived well by the reviewers. 
        \item The claims made should match theoretical and experimental results, and reflect how much the results can be expected to generalize to other settings. 
        \item It is fine to include aspirational goals as motivation as long as it is clear that these goals are not attained by the paper. 
    \end{itemize}

\item {\bf Limitations}
    \item[] Question: Does the paper discuss the limitations of the work performed by the authors?
    \item[] Answer: \answerYes{}
    \item[] Justification: Section~\ref{sec:disc} (Discussion and limitations) lists cap selection bias and transfer, task independence, fixed rivals, noise, rule coverage, snapshots and unfinished items.
    \item[] Guidelines:
    \begin{itemize}
        \item The answer \answerNA{} means that the paper has no limitation while the answer \answerNo{} means that the paper has limitations, but those are not discussed in the paper. 
        \item The authors are encouraged to create a separate ``Limitations'' section in their paper.
        \item The paper should point out any strong assumptions and how robust the results are to violations of these assumptions (e.g., independence assumptions, noiseless settings, model well-specification, asymptotic approximations only holding locally). The authors should reflect on how these assumptions might be violated in practice and what the implications would be.
        \item The authors should reflect on the scope of the claims made, e.g., if the approach was only tested on a few datasets or with a few runs. In general, empirical results often depend on implicit assumptions, which should be articulated.
        \item The authors should reflect on the factors that influence the performance of the approach. For example, a facial recognition algorithm may perform poorly when image resolution is low or images are taken in low lighting. Or a speech-to-text system might not be used reliably to provide closed captions for online lectures because it fails to handle technical jargon.
        \item The authors should discuss the computational efficiency of the proposed algorithms and how they scale with dataset size.
        \item If applicable, the authors should discuss possible limitations of their approach to address problems of privacy and fairness.
        \item While the authors might fear that complete honesty about limitations might be used by reviewers as grounds for rejection, a worse outcome might be that reviewers discover limitations that aren't acknowledged in the paper. The authors should use their best judgment and recognize that individual actions in favor of transparency play an important role in developing norms that preserve the integrity of the community. Reviewers will be specifically instructed to not penalize honesty concerning limitations.
    \end{itemize}

\item {\bf Theory assumptions and proofs}
    \item[] Question: For each theoretical result, does the paper provide the full set of assumptions and a complete (and correct) proof?
    \item[] Answer: \answerYes{}
    \item[] Justification: Propositions~\ref{prop:mean}--\ref{prop:median} state their assumptions in Section~\ref{sec:solvers} and are proved in Appendix~\ref{app:proofs}; rank-rule results are computational and no theorem is claimed.
    \item[] Guidelines:
    \begin{itemize}
        \item The answer \answerNA{} means that the paper does not include theoretical results. 
        \item All the theorems, formulas, and proofs in the paper should be numbered and cross-referenced.
        \item All assumptions should be clearly stated or referenced in the statement of any theorems.
        \item The proofs can either appear in the main paper or the supplemental material, but if they appear in the supplemental material, the authors are encouraged to provide a short proof sketch to provide intuition. 
        \item Inversely, any informal proof provided in the core of the paper should be complemented by formal proofs provided in appendix or supplemental material.
        \item Theorems and Lemmas that the proof relies upon should be properly referenced. 
    \end{itemize}

    \item {\bf Experimental result reproducibility}
    \item[] Question: Does the paper fully disclose all the information needed to reproduce the main experimental results of the paper to the extent that it affects the main claims and/or conclusions of the paper (regardless of whether the code and data are provided or not)?
    \item[] Answer: \answerYes{}
    \item[] Justification: Section~\ref{sec:setup}--\ref{sec:exp} define rules, caps, margin and filters; \texttt{REPRODUCE.md} in the supplementary code gives data versions, seeds and commands.
    \item[] Guidelines:
    \begin{itemize}
        \item The answer \answerNA{} means that the paper does not include experiments.
        \item If the paper includes experiments, a \answerNo{} answer to this question will not be perceived well by the reviewers: Making the paper reproducible is important, regardless of whether the code and data are provided or not.
        \item If the contribution is a dataset and\slash or model, the authors should describe the steps taken to make their results reproducible or verifiable. 
        \item Depending on the contribution, reproducibility can be accomplished in various ways. For example, if the contribution is a novel architecture, describing the architecture fully might suffice, or if the contribution is a specific model and empirical evaluation, it may be necessary to either make it possible for others to replicate the model with the same dataset, or provide access to the model. In general. releasing code and data is often one good way to accomplish this, but reproducibility can also be provided via detailed instructions for how to replicate the results, access to a hosted model (e.g., in the case of a large language model), releasing of a model checkpoint, or other means that are appropriate to the research performed.
        \item While NeurIPS does not require releasing code, the conference does require all submissions to provide some reasonable avenue for reproducibility, which may depend on the nature of the contribution. For example
        \begin{enumerate}
            \item If the contribution is primarily a new algorithm, the paper should make it clear how to reproduce that algorithm.
            \item If the contribution is primarily a new model architecture, the paper should describe the architecture clearly and fully.
            \item If the contribution is a new model (e.g., a large language model), then there should either be a way to access this model for reproducing the results or a way to reproduce the model (e.g., with an open-source dataset or instructions for how to construct the dataset).
            \item We recognize that reproducibility may be tricky in some cases, in which case authors are welcome to describe the particular way they provide for reproducibility. In the case of closed-source models, it may be that access to the model is limited in some way (e.g., to registered users), but it should be possible for other researchers to have some path to reproducing or verifying the results.
        \end{enumerate}
    \end{itemize}


\item {\bf Open access to data and code}
    \item[] Question: Does the paper provide open access to the data and code, with sufficient instructions to faithfully reproduce the main experimental results, as described in supplemental material?
    \item[] Answer: \answerYes{}
    \item[] Justification: The data are public (Open LLM Leaderboard and HELM, versions and checksums in \texttt{REPRODUCE.md}); the supplementary code reproduces every table and figure with the commands listed there.
    \item[] Guidelines:
    \begin{itemize}
        \item The answer \answerNA{} means that paper does not include experiments requiring code.
        \item Please see the NeurIPS code and data submission guidelines (\url{https://neurips.cc/public/guides/CodeSubmissionPolicy}) for more details.
        \item While we encourage the release of code and data, we understand that this might not be possible, so \answerNo{} is an acceptable answer. Papers cannot be rejected simply for not including code, unless this is central to the contribution (e.g., for a new open-source benchmark).
        \item The instructions should contain the exact command and environment needed to run to reproduce the results. See the NeurIPS code and data submission guidelines (\url{https://neurips.cc/public/guides/CodeSubmissionPolicy}) for more details.
        \item The authors should provide instructions on data access and preparation, including how to access the raw data, preprocessed data, intermediate data, and generated data, etc.
        \item The authors should provide scripts to reproduce all experimental results for the new proposed method and baselines. If only a subset of experiments are reproducible, they should state which ones are omitted from the script and why.
        \item At submission time, to preserve anonymity, the authors should release anonymized versions (if applicable).
        \item Providing as much information as possible in supplemental material (appended to the paper) is recommended, but including URLs to data and code is permitted.
    \end{itemize}


\item {\bf Experimental setting/details}
    \item[] Question: Does the paper specify all the training and test details (e.g., data splits, hyperparameters, how they were chosen, type of optimizer) necessary to understand the results?
    \item[] Answer: \answerYes{}
    \item[] Justification: Section~\ref{sec:exp} and Appendix state data filters, dedupe, caps, bootstrap and Monte Carlo settings; \texttt{REPRODUCE.md} lists the rest.
    \item[] Guidelines:
    \begin{itemize}
        \item The answer \answerNA{} means that the paper does not include experiments.
        \item The experimental setting should be presented in the core of the paper to a level of detail that is necessary to appreciate the results and make sense of them.
        \item The full details can be provided either with the code, in appendix, or as supplemental material.
    \end{itemize}

\item {\bf Experiment statistical significance}
    \item[] Question: Does the paper report error bars suitably and correctly defined or other appropriate information about the statistical significance of the experiments?
    \item[] Answer: \answerYes{}
    \item[] Justification: Bootstrap intervals over targets (Figure~\ref{fig:main}, Tables~\ref{tab:T3}, \ref{tab:T4}) and percentile ranges over Monte Carlo draws for evaluation noise and task resampling on OLLB and HELM (Figure~\ref{fig:robust}, Section~\ref{sec:results}); these cover seven OLLB and six HELM targets, not all targets, and are stated as spreads of draws, not confidence intervals.
    \item[] Guidelines:
    \begin{itemize}
        \item The answer \answerNA{} means that the paper does not include experiments.
        \item The authors should answer \answerYes{} if the results are accompanied by error bars, confidence intervals, or statistical significance tests, at least for the experiments that support the main claims of the paper.
        \item The factors of variability that the error bars are capturing should be clearly stated (for example, train/test split, initialization, random drawing of some parameter, or overall run with given experimental conditions).
        \item The method for calculating the error bars should be explained (closed form formula, call to a library function, bootstrap, etc.)
        \item The assumptions made should be given (e.g., Normally distributed errors).
        \item It should be clear whether the error bar is the standard deviation or the standard error of the mean.
        \item It is OK to report 1-sigma error bars, but one should state it. The authors should preferably report a 2-sigma error bar than state that they have a 96\% CI, if the hypothesis of Normality of errors is not verified.
        \item For asymmetric distributions, the authors should be careful not to show in tables or figures symmetric error bars that would yield results that are out of range (e.g., negative error rates).
        \item If error bars are reported in tables or plots, the authors should explain in the text how they were calculated and reference the corresponding figures or tables in the text.
    \end{itemize}

\item {\bf Experiments compute resources}
    \item[] Question: For each experiment, does the paper provide sufficient information on the computer resources (type of compute workers, memory, time of execution) needed to reproduce the experiments?
    \item[] Answer: \answerYes{}
    \item[] Justification: Section~\ref{sec:repro} and \texttt{REPRODUCE.md} give the CPU type, memory and approximate wall-clock per step on a shared container; CPU only, no GPU.
    \item[] Guidelines:
    \begin{itemize}
        \item The answer \answerNA{} means that the paper does not include experiments.
        \item The paper should indicate the type of compute workers CPU or GPU, internal cluster, or cloud provider, including relevant memory and storage.
        \item The paper should provide the amount of compute required for each of the individual experimental runs as well as estimate the total compute. 
        \item The paper should disclose whether the full research project required more compute than the experiments reported in the paper (e.g., preliminary or failed experiments that didn't make it into the paper). 
    \end{itemize}
    
\item {\bf Code of ethics}
    \item[] Question: Does the research conducted in the paper conform, in every respect, with the NeurIPS Code of Ethics \url{https://neurips.cc/public/EthicsGuidelines}?
    \item[] Answer: \answerYes{}
    \item[] Justification: The work uses only public score tables, trains no models and involves no human subjects; we have read the NeurIPS Code of Ethics and see no conflict.
    \item[] Guidelines:
    \begin{itemize}
        \item The answer \answerNA{} means that the authors have not reviewed the NeurIPS Code of Ethics.
        \item If the authors answer \answerNo, they should explain the special circumstances that require a deviation from the Code of Ethics.
        \item The authors should make sure to preserve anonymity (e.g., if there is a special consideration due to laws or regulations in their jurisdiction).
    \end{itemize}


\item {\bf Broader impacts}
    \item[] Question: Does the paper discuss both potential positive societal impacts and negative societal impacts of the work performed?
    \item[] Answer: \answerYes{}
    \item[] Justification: Section~\ref{sec:disc} (Broader impact) discusses benefits for leaderboard designers and the risk that the analysis is read as a manipulation recipe.
    \item[] Guidelines:
    \begin{itemize}
        \item The answer \answerNA{} means that there is no societal impact of the work performed.
        \item If the authors answer \answerNA{} or \answerNo, they should explain why their work has no societal impact or why the paper does not address societal impact.
        \item Examples of negative societal impacts include potential malicious or unintended uses (e.g., disinformation, generating fake profiles, surveillance), fairness considerations (e.g., deployment of technologies that could make decisions that unfairly impact specific groups), privacy considerations, and security considerations.
        \item The conference expects that many papers will be foundational research and not tied to particular applications, let alone deployments. However, if there is a direct path to any negative applications, the authors should point it out. For example, it is legitimate to point out that an improvement in the quality of generative models could be used to generate Deepfakes for disinformation. On the other hand, it is not needed to point out that a generic algorithm for optimizing neural networks could enable people to train models that generate Deepfakes faster.
        \item The authors should consider possible harms that could arise when the technology is being used as intended and functioning correctly, harms that could arise when the technology is being used as intended but gives incorrect results, and harms following from (intentional or unintentional) misuse of the technology.
        \item If there are negative societal impacts, the authors could also discuss possible mitigation strategies (e.g., gated release of models, providing defenses in addition to attacks, mechanisms for monitoring misuse, mechanisms to monitor how a system learns from feedback over time, improving the efficiency and accessibility of ML).
    \end{itemize}
    
\item {\bf Safeguards}
    \item[] Question: Does the paper describe safeguards that have been put in place for responsible release of data or models that have a high risk for misuse (e.g., pre-trained language models, image generators, or scraped datasets)?
    \item[] Answer: \answerNA{}
    \item[] Justification: No models or data with high misuse risk are released; only derived score tables and code.
    \item[] Guidelines:
    \begin{itemize}
        \item The answer \answerNA{} means that the paper poses no such risks.
        \item Released models that have a high risk for misuse or dual-use should be released with necessary safeguards to allow for controlled use of the model, for example by requiring that users adhere to usage guidelines or restrictions to access the model or implementing safety filters. 
        \item Datasets that have been scraped from the Internet could pose safety risks. The authors should describe how they avoided releasing unsafe images.
        \item We recognize that providing effective safeguards is challenging, and many papers do not require this, but we encourage authors to take this into account and make a best faith effort.
    \end{itemize}

\item {\bf Licenses for existing assets}
    \item[] Question: Are the creators or original owners of assets (e.g., code, data, models), used in the paper, properly credited and are the license and terms of use explicitly mentioned and properly respected?
    \item[] Answer: \answerYes{}
    \item[] Justification: Sources and licence status are listed in \texttt{REPRODUCE.md}: neither leaderboard declares a licence for its score data, HELM code is Apache-2.0; we say so and release only derived tables.
    \item[] Guidelines:
    \begin{itemize}
        \item The answer \answerNA{} means that the paper does not use existing assets.
        \item The authors should cite the original paper that produced the code package or dataset.
        \item The authors should state which version of the asset is used and, if possible, include a URL.
        \item The name of the license (e.g., CC-BY 4.0) should be included for each asset.
        \item For scraped data from a particular source (e.g., website), the copyright and terms of service of that source should be provided.
        \item If assets are released, the license, copyright information, and terms of use in the package should be provided. For popular datasets, \url{paperswithcode.com/datasets} has curated licenses for some datasets. Their licensing guide can help determine the license of a dataset.
        \item For existing datasets that are re-packaged, both the original license and the license of the derived asset (if it has changed) should be provided.
        \item If this information is not available online, the authors are encouraged to reach out to the asset's creators.
    \end{itemize}

\item {\bf New assets}
    \item[] Question: Are new assets introduced in the paper well documented and is the documentation provided alongside the assets?
    \item[] Answer: \answerYes{}
    \item[] Justification: The released code and derived score tables are documented in \texttt{REPRODUCE.md}; no other new assets are released.
    \item[] Guidelines:
    \begin{itemize}
        \item The answer \answerNA{} means that the paper does not release new assets.
        \item Researchers should communicate the details of the dataset\slash code\slash model as part of their submissions via structured templates. This includes details about training, license, limitations, etc. 
        \item The paper should discuss whether and how consent was obtained from people whose asset is used.
        \item At submission time, remember to anonymize your assets (if applicable). You can either create an anonymized URL or include an anonymized zip file.
    \end{itemize}

\item {\bf Crowdsourcing and research with human subjects}
    \item[] Question: For crowdsourcing experiments and research with human subjects, does the paper include the full text of instructions given to participants and screenshots, if applicable, as well as details about compensation (if any)? 
    \item[] Answer: \answerNA{}
    \item[] Justification: The paper involves no crowdsourcing or human subjects.
    \item[] Guidelines:
    \begin{itemize}
        \item The answer \answerNA{} means that the paper does not involve crowdsourcing nor research with human subjects.
        \item Including this information in the supplemental material is fine, but if the main contribution of the paper involves human subjects, then as much detail as possible should be included in the main paper. 
        \item According to the NeurIPS Code of Ethics, workers involved in data collection, curation, or other labor should be paid at least the minimum wage in the country of the data collector. 
    \end{itemize}

\item {\bf Institutional review board (IRB) approvals or equivalent for research with human subjects}
    \item[] Question: Does the paper describe potential risks incurred by study participants, whether such risks were disclosed to the subjects, and whether Institutional Review Board (IRB) approvals (or an equivalent approval/review based on the requirements of your country or institution) were obtained?
    \item[] Answer: \answerNA{}
    \item[] Justification: The paper involves no human subjects, so no IRB approval applies.
    \item[] Guidelines:
    \begin{itemize}
        \item The answer \answerNA{} means that the paper does not involve crowdsourcing nor research with human subjects.
        \item Depending on the country in which research is conducted, IRB approval (or equivalent) may be required for any human subjects research. If you obtained IRB approval, you should clearly state this in the paper. 
        \item We recognize that the procedures for this may vary significantly between institutions and locations, and we expect authors to adhere to the NeurIPS Code of Ethics and the guidelines for their institution. 
        \item For initial submissions, do not include any information that would break anonymity (if applicable), such as the institution conducting the review.
    \end{itemize}

\item {\bf Declaration of LLM usage}
    \item[] Question: Does the paper describe the usage of LLMs if it is an important, original, or non-standard component of the core methods in this research? Note that if the LLM is used only for writing, editing, or formatting purposes and does \emph{not} impact the core methodology, scientific rigor, or originality of the research, declaration is not required.
    %this research? 
    \item[] Answer: \answerNA{}
    \item[] Justification: No LLM is part of the method; LLM assistance was used only for writing and coding assistance.
    \item[] Guidelines:
    \begin{itemize}
        \item The answer \answerNA{} means that the core method development in this research does not involve LLMs as any important, original, or non-standard components.
        \item Please refer to our LLM policy in the NeurIPS handbook for what should or should not be described.
    \end{itemize}

\end{enumerate}

\end{document}
